\question
Идёт практика в вузе по предмету «Математический анализ». На практике присутствует всего 5 студентов из группы. Умник A, его друг B, а также студенты C, D, E. Студентов в группе так мало, потому что все остальные отчислены. Причиной такого массового отчисления стала следующая закономерность, выявленная на этих практиках за последнее время:

\begin{enumerate}
  \item Если A может решить задачу, то и B может её решить.
  \item Хотя бы один из D или E может решить задачу.
  \item Либо B, либо C решали такую задачу ранее и могут решить её снова.
  \item C и D смогут решить задачу только если они совместно ботали, либо совместно не ботали.
  \item Если E решает задачу, то A и D понимают его логику и могут решить эту задачу.
\end{enumerate}
Определите, кто из студентов гарантированно может решить задачу на текущей паре. Для каких студентов приведённых сведений недостаточно, чтобы установить, могут ли они решить задачу? Обоснуйте ответ.

---------------

Автор -- Вероника Вдовина